\documentclass[11pt]{article}

\usepackage[margin=0.95in]{geometry}
\usepackage[T1]{fontenc}
\usepackage{lmodern}
\usepackage[numbers,sort&compress]{natbib}
\usepackage{booktabs,tabularx,array}
\usepackage{microtype}
\usepackage{xcolor}
\usepackage{hyperref}
\usepackage{fancyhdr}
\usepackage{parskip}

\definecolor{icpnavy}{HTML}{183B56}
\definecolor{icpteal}{HTML}{2B6F77}
\definecolor{icpgold}{HTML}{A56A18}
\definecolor{icpgray}{HTML}{5E6870}
\definecolor{icppale}{HTML}{F1F5F6}
\definecolor{icpwarm}{HTML}{FBF5EA}

\hypersetup{
  colorlinks=true,
  linkcolor=icpnavy,
  citecolor=icpteal,
  urlcolor=icpteal,
  pdftitle={How the Past Changes What Can Happen Next},
  pdfauthor={Douglas Clark}
}

\pagestyle{fancy}
\fancyhf{}
\fancyhead[L]{\footnotesize\color{icpgray}A plain-language companion to the IOTA/ICP scientific synthesis}
\fancyhead[R]{\footnotesize\color{icpgray}July 2026}
\fancyfoot[C]{\thepage}
\renewcommand{\headrulewidth}{0.3pt}
\setlength{\headheight}{14pt}
\setlength{\parindent}{0pt}
\setlength{\parskip}{0.72em}
\setlength{\leftmargini}{1.7em}
\emergencystretch=2em

\newcommand{\plainbox}[2]{%
  \begin{center}
  \colorbox{#1}{\parbox{0.92\linewidth}{\vspace{0.35em}#2\vspace{0.35em}}}
  \end{center}}

\title{\vspace{-1.2em}\textbf{How the Past Changes What Can Happen Next}\\
\large A true-layperson guide to IOTA and the Information--Constraint Principle}
\author{Douglas Clark}
\date{Plain-language companion --- 20 July 2026}

\begin{document}
\maketitle

\plainbox{icppale}{\small\textbf{The answer in one paragraph.}
IOTA and the Information--Constraint Principle are worth pursuing as a
scientific research program. The project has already shown, inside carefully
defined computer models, that a useful relational role can disappear in one
place and reappear through a completely different local arrangement. Other
sciences show many parts of the larger idea: memories survive their initiating
events, physical records alter later behavior, and organized quantum systems
can produce particle-like patterns. But no current result shows that the
universe is made of an IOTA lattice, that information is the unique foundation
of nature, or that matter and energy have been derived from it. The idea has
legs; it has not yet crossed the finish line.}

\section*{The question that started the journey}

How can the effect of something survive after the original thing is gone?

A foot lifts from wet ground, but the footprint remains. A bell stops moving,
but a recording can reproduce its sound. A virus attacks a bacterium; much
later, descendants of that bacterium can carry a molecular record of the
attack. In each case, something about an earlier encounter matters later even
though the original actor or arrangement is no longer present.

This can sound mysterious only because the word \emph{information} is used for
several different things. The footprint is not the foot. It is also not a
ghostly message floating outside the physical world. The mud has been changed.
The recording is not the bell, but a new physical arrangement carries a usable
pattern. The bacterial record is not the virus, but molecules in the cell now
make a different response possible.

That distinction is the starting point for IOTA:

\plainbox{icpwarm}{\small
Information may persist \emph{beyond its original carrier}. That does not mean
it persists \emph{without any physical carrier}.}

The deeper question is whether this familiar kind of persistence is a clue to
something more basic. Could reality be built, at its lowest level, from events,
relationships, and rules about which events may follow? Could the things we
usually treat as fundamental---particles, fields, mass, energy, space, and
time---instead be durable patterns that emerge from that deeper organization?

IOTA explores that possibility. The Information--Constraint Principle, called
the principle in the rest of this essay, is the proposed bridge between the
basic layer and the world we know. This paper asks two questions. Does that
proposal fit what science currently knows well enough to deserve serious work?
And if it does, what must happen next?

\section*{The idea in ordinary language}

The project's working definition is deliberately practical:

\plainbox{icppale}{\small\textbf{IOTA's premise.}
Information is a change in what can happen next, caused by an interaction.}

Imagine a railway yard. The same trains can reach different destinations
depending on how the switches are set. The switches do not push every train to
one inevitable destination. They shape the routes that are open, closed, easy,
or unlikely. An earlier worker may set a switch and leave. The worker is gone,
but the setting still changes the journeys available to later trains.

This is what \emph{constraint} means here. It is not merely a restriction and
it is not a new invisible force. A constraint can close a path, open a path,
redirect traffic, or change the odds among paths. A bridge can make a
destination reachable. A memory can make one response more likely than
another. An arrangement can preserve useful options at one level while ruling
out movements at another.

The phrase \emph{what can happen next} also needs care. It never means every
future anyone can imagine. It means the futures available under stated rules,
conditions, actions, measuring methods, resources, and time span. A locked
door is impassable to a person with bare hands, but not to a firefighter with
an axe. Both statements can be true because they describe different allowed
actions. Mathematics for reachable and survivable futures already exists in
control and viability theory \citep{SaintPierre1994Viability,MitchellEtAl2005Reachable}.
Any new principle must respect and build beyond that work.

IOTA makes a larger proposal about the foundation of physics. Its candidate
starting layer contains local distinctions, relationships, and rules of
succession. The present computer lattice is one concrete way to test the idea.
It is \emph{not} being declared the literal grid of the cosmos. If the research
program succeeds, familiar physics would have to emerge in stages:

\begin{enumerate}
  \item interactions produce differences in later possibilities;
  \item some differences persist and can pass to new physical carriers;
  \item some become durable, coordinated patterns;
  \item some patterns acquire particle-like behavior;
  \item interacting patterns recover fields and conserved quantities; and
  \item mass, energy, geometry, clocks, and ordinary large-scale physics emerge
  as reliable features of that organization.
\end{enumerate}

This is a research ladder, not a list of completed discoveries. The present
project is near its first rungs.

\section*{What the IOTA project has actually shown}

The cleanest project result concerns a bridge. In the computer model, two
regions were connected through a local arrangement of cells. That original
bridge disappeared. Later, the regions became connected again through a new
bridge that shared no cells with the first one. Removing the new bridge then
destroyed the renewed connection. The new arrangement was therefore doing the
work; it was not merely near the place where the work happened
\citep{Clark2026IotaConstraint}.

The numbers matter. Thirty-six candidate relocation events were replayed.
Twelve had the right kind of original bridge, and three events across two
independent starting cases passed the strict replacement test. That is a real
positive result, but a narrow one. It shows that this kind of role replacement
\emph{can} happen in the declared model. It does not show that it always
happens, lasts forever, or occurs in nature.

Follow-up tests made the picture richer. Some relationships survived after the
original cell addresses died because the project tracked a continuing line of
expression rather than a fixed address. Some lasted much longer than others.
Opening a blocked route did not guarantee that the route would eventually be
used: timing and competing events could close the opportunity first. When an
observed successor was prevented, the model often found another successor.
This exposed hidden backups. A role could be necessary only after several
alternatives had been removed in sequence.

These are useful lessons about possibility. An open route is not the same as a
realized outcome. A thing that looks necessary may have an unseen substitute.
A record tied to one address may look dead even while its role continues
elsewhere. The right object may be a continuing relation, not a marked lump
moving across the screen.

The project also kept its failures. One proposed way of describing a graph
role did not pass its full confirmation test. A later independent repair study
did not collect the minimum number of eligible cases, so its outcome was not
opened and the idea remains untested. These are not embarrassing details. They
are evidence that the tests can say no.

A second project asked whether two local quantum models could reproduce a
limited part of a well-known equation used for electron-like particles. They
did pass the surviving tests for that limited role \citep{Clark2026DiracRole}.
But the construction began with quantum rules already supplied. It also
supplied the clock, the rule that turns a quantum description into measurement
probabilities, and a mass setting. An audit found a faulty test that had looked
at an empty part of the data; those results were withdrawn. The audit also
showed that the two free models were more closely related than first thought.

The honest conclusion is encouraging and restrained: local mechanisms can
reproduce a bounded piece of known quantum behavior, and the audit improved
the claim. This is not yet a derivation of quantum theory, a particle, or mass
from base IOTA.

A third project studied the geometry of meaning in an engineered language
system \citep{Clark2026PSA}. It showed that relations to fixed reference
statements can preserve useful structure and that comparisons made from
separate reference sets can be checked against one another. It also used
negative controls to show what the method could not establish. This is useful
evidence about relational measurement. It is not evidence that nature is made
of language or that an artificial-intelligence embedding is fundamental
physics.

\section*{What the wider world already teaches us}

The larger scientific literature does not prove IOTA. It does show that the
program is asking a coherent question and that many pieces of the proposed
mechanism are physically possible.

\subsection*{Records can leave their source}

In quantum experiments, a system can interact with its surroundings so that
different pieces of the surroundings carry readable traces of the system.
Under the right conditions, separate observers can learn the same limited
fact by measuring different environmental pieces rather than touching the
source. This area is sometimes called quantum Darwinism. Its lesson is modest
but important: an interaction can create secondary physical carriers of a
record
\citep{OllivierPoulinZurek2004Witness,BrandaoEtAl2015Darwinism,CiampiniEtAl2018DarwinismExperiment,ChenEtAl2019DarwinismExperiment,UndenEtAl2019Classicality}.

It does not follow that every environment makes good records, that the record
exists without a carrier, or that there is always one uniquely correct way to
divide the world into ``system'' and ``environment.'' A correlation is not
automatically an accessible memory, and a memory is not automatically a new
law of nature. Modern descriptions of multi-step quantum experiments already
capture how earlier interactions can influence later responses
\citep{PollockEtAl2018ProcessTensor,PollockEtAl2018Markov}. The principle must
add something testable beyond a complete version of that account.

\subsection*{Living things are made of maintained history}

Biology supplies especially vivid examples. After certain bacteria survive a
virus, they can insert a short piece of viral sequence into their own genetic
record. Editing that record changes which virus they resist later
\citep{BarrangouEtAl2007CRISPR}. In yeast,
an earlier period of nutrient shortage can establish a chromatin state that
outlasts the factor that initiated it and survives several cell divisions.
Remove the machinery that maintains the mark, however, and the memory fades
\citep{LightEtAl2010INO1,SumpEtAl2022INO1Memory}.

Evolution also has a path. In a classic protein experiment, researchers built
every combination of five mutations. Under the experiment's stated conditions,
only a small fraction of the direct step-by-step routes kept improving at every
step. Earlier changes made some later changes workable and others damaging.
The result does not mean that nature had only those routes under every
population, environment, and time span. It means that history changed
accessibility under the declared test
\citep{WeinreichEtAl2006MutationalPaths,HarmsThornton2014Contingency}.

Experiments with self-copying molecules and cell-like compartments tell a
similar story. Cooperative networks can amplify one another. Boundaries can
protect useful replicators from parasites. Yet self-copying, self-maintenance,
heredity, and open-ended evolution are not the same achievement. Many synthetic
systems still depend on researchers supplying food, enzymes, timing, or
division. Biology therefore supports physically maintained, history-dependent
constraint. It does not establish a new foundation of physics
\citep{LincolnJoyce2009RNAReplication,VasasEtAl2010Evolvability,VaidyaEtAl2012CooperativeRNA,IchihashiEtAl2013RNACompartment}.

\subsection*{Information can change the work a machine can obtain}

A controller that records where a tiny moving particle is can use that record
to choose a later action and gain free energy---energy available to do useful
work. This has been shown in a laboratory experiment and fits established
thermodynamic accounting
\citep{SagawaUeda2010Feedback,ToyabeEtAl2010InformationEnergy}. The information
is not free magic: the measuring device, memory, controller, reset procedure,
heat bath, and allowed actions all matter.

A 2026 experimental preprint gives a close modern example. A visible particle
in a complex fluid appeared to remember its history because hidden parts of
the surrounding fluid relaxed slowly. Using delayed measurements, researchers
reported extracting work associated with that hidden memory
\citep{MurugaEtAl2026HiddenDegrees}. This is exciting precisely because it
makes the boundary clear. The memory is outside the currently observed
particle, but it remains physically carried by the environment. Once those
hidden parts are included, ordinary physics may fully explain the effect.

\subsection*{Particle-like things can emerge from organized many-part systems}

The word \emph{particle} can suggest a tiny hard bead. Physics is subtler.
Inside materials and engineered quantum systems, coordinated disturbances can
behave like particles: they can travel, carry charge-like properties, collide,
or obey unusual exchange rules. Physicists call these collective objects
\emph{quasiparticles}.

A stadium wave is a useful first picture, as long as we do not mistake it for a
quantum theory. The wave moves around the stadium while individual people only
stand and sit. The moving pattern has a life and behavior not captured by
following one person. In real quantum systems, the collective behavior is far
richer and is backed by exact models and controlled experiments. Local quantum
models can produce fractional and particle-like excitations
\citep{LevinWen2005StringNet,Kitaev2006Anyons}; quantum processors and atomic
gases have engineered and measured striking examples
\citep{IqbalEtAl2024NonAbelian,DharEtAl2025Anyonization}.

This defeats a blanket objection: particle-like behavior \emph{can} emerge
from organized local ingredients. It does not prove that IOTA's present
patterns are particles. Every example just named already assumes quantum
mechanics. A cellular quantum model can approach a limited free-field theory
in one dimension, while related constructions face serious barriers in higher
dimensions \citep{MlodinowBrun2020QCAQFT}. IOTA must build the quantum layer it
currently borrows.

\section*{Three tempting clues that are not proof}

Some of the most exciting subjects in science are useful to this project
because they expose mistakes the project must avoid.

\subsection*{Zipf's law: the same shape can have many causes}

Put words in order from most common to least common and their frequencies often
fall in a remarkably regular way. Similar rank patterns appear in other kinds
of data. This is Zipf's law. It is tempting to see such a clean curve as the
signature of a deep organizing principle.

But very different mechanisms can make Zipf-like patterns: random text,
growth with preferential reuse, hidden variables, least-effort models, and
history-dependent shrinking of the available choices all produce related
shapes under different conditions
\citep{Simon1955Skew,Li1992RandomTexts,FerrerCanchoSole2003LeastEffort,SchwabEtAl2014LatentVariables,CorominasMurtraEtAl2015SSR}.
Even deciding whether data truly follow a power law requires careful comparison
with rival distributions \citep{ClausetEtAl2009PowerLaws}.

So a Zipf curve cannot prove IOTA. At best it can suggest a question. The useful
test would hold the visible rank curve nearly the same while changing the
history or the proposed carrier, then ask whether later responses differ. If
two systems share the curve but not their future behavior, the pattern alone
was never the mechanism.

\subsection*{Entropy: one number can hide different futures}

Entropy is not one universal synonym for disorder. In communication, it can
measure uncertainty in a set of messages \citep{Shannon1948Communication}. In
thermodynamics, it is tied to physical states, heat, work, and the direction of
large-scale change \citep{Boltzmann1877Relationship,LiebYngvason1999SecondLaw}.
Other entropies describe dynamical unpredictability, quantum resources, or
coarse-grained descriptions.

Constraint and entropy are related, but neither simply replaces the other.
Closing paths can lower uncertainty in one description. Maintaining a living
cell can keep a narrow set of local conditions by releasing heat elsewhere.
Two systems can have the same entropy value while allowing different outcomes,
remembering different histories, or requiring different resources to reverse.
The answer changes when we change which states count as distinct, what can be
measured, what actions are allowed, and how long we wait.

For the principle, entropy is therefore a check, not a shortcut. Every claim
must state the menu of states, the probabilities or physical measure, the
boundary, the allowed actions, and the time span. A claim that merely renames
low entropy as ``more constraint'' would add little. More promising is to ask
whether an embodied record changes the structure of future paths even when a
single entropy score is held fixed. Work connecting prediction and wasted
energy shows why stored history can matter without making information a new
substance \citep{StillEtAl2012Prediction}.

\subsection*{Black holes: preservation is not the same as recovery}

Black holes make the persistence question unavoidable. If an object falls in
and the black hole later evaporates, what happens to the distinctions that once
described the object? Modern theory has made major progress. Calculations in
controlled gravitational models can reproduce the expected rise and fall of
the radiation's entropy, often called a Page curve
\citep{Page1993Information,AlmheiriEtAl2019BulkEntropy,PeningtonEtAl2022ReplicaWormholes}.
Ideas from quantum error correction also explain how information about one
region can be represented in a distributed and indirect way
\citep{AlmheiriDongHarlow2015QEC}.

But several questions must not be collapsed into one. Is information preserved
in the complete state? Can a chosen set of measurements distinguish it? Can it
be reconstructed in principle? Can an actual observer recover a message? Can
that recovery be done before the required time or machinery becomes absurd?
Computational barriers can separate ``present in principle'' from ``usable''
\citep{HarlowHayden2013Computation}.

Page curves, islands, and reconstruction maps are not identified physical
carriers in the simple sense of a memory chip or a molecule. The strongest
results are theoretical and often belong to carefully controlled models, not
to observed evaporation of an astrophysical black hole. Black-hole information
theory is therefore an exceptional stress test for the principle's language of
carriers, access, boundaries, and recovery. It is not evidence that IOTA is the
substrate of the universe.

\section*{The strongest skeptical answer}

A critic can make a powerful and simple objection:

\plainbox{icpwarm}{\small
If we know the complete physical state, the laws, the allowed actions, and the
boundary, then ordinary physics already tells us the possible futures.
Calling part of that state ``information'' and its effects ``constraint'' may
only rename what was already there.}

This objection is not a misunderstanding. It is the main scientific test.
Control theory already calculates reachability. Predictive models already
group histories that lead to the same future probabilities. Causal methods
already ask what changes after an intervention. One established framework even
scrambles correlations and measures whether an organism-like system remains
viable \citep{KolchinskyWolpert2018Semantic}. The broad claim ``information
changes future possibilities'' is not new.

For the proposed principle to become more than a useful viewpoint, carrier
succession must earn at least one concrete payoff. It might reveal a smallest
physical description that predicts all relevant responses. It might identify
a feature that remains stable across very different physical realizations. It
might prove that some transformation is impossible. It might compress many
intervention results better than existing methods. Best of all, it might make
a new prediction that ordinary complete-state accounts do not make.

If none of those happens, IOTA could still offer a productive vocabulary and a
good way to design experiments. It would not be a distinct law of physics.

\section*{The honest verdict}

The verdict is positive, but graded:

\plainbox{icppale}{\small\textbf{IOTA and the principle have real scientific
legs as a research program and an operational way to connect persistent
records with future possibilities. They are not established as a fundamental
law or as the foundation of nature.}}

That conclusion rests on more than analogy. The central question can be tested.
The project has produced bounded causal demonstrations, including replacement
of an original bridge by a disjoint successor. Independent sciences establish
that records can be physically copied, maintained, moved, and used; that
history can reshape biological and thermodynamic responses; and that local
quantum organization can support particle-like roles. There is also a clear
sequence of experiments that could strengthen or defeat the larger proposal.

The ceiling on the conclusion matters just as much. The positive IOTA results
come from one family of computer rules and representations. The quantum work
starts with quantum mechanics rather than deriving it. No external experiment
has tested the full carrier-replacement criterion as an IOTA prediction. No
result derives mass, energy, space, time, or the known laws from an
informational substrate.

\subsection*{What we know}

\begin{itemize}
  \item Earlier interactions can leave physical records that affect later
  responses.
  \item Such records can outlast the event, object, or local arrangement that
  first created them while remaining carried by something physical.
  \item Reachable futures depend on rules, state, boundary, allowed actions,
  resources, and time span.
  \item Collective quantum organization can produce genuine particle-like
  behavior.
  \item The IOTA model has produced limited, intervention-tested examples of a
  relational role continuing through local carrier replacement.
\end{itemize}

\subsection*{What we infer}

\begin{itemize}
  \item Formation, persistence, and transfer of constraint may be a useful
  common language across physics, biology, and complex systems.
  \item Carrier succession may be a better target than searching for a tiny
  object that literally transports a payload.
  \item A ladder from persistent relations to coherent collective objects is
  scientifically plausible enough to attempt.
\end{itemize}

These are synthesis judgments. They are not results proved by any one cited
field.

\subsection*{What we do not know}

\begin{itemize}
  \item Whether information, rather than matter, energy, quantum fields, or
  some other structure, is the most basic layer of nature.
  \item Whether the current IOTA rules describe anything beyond their computer
  model.
  \item Whether IOTA can generate quantum states and measurement probabilities
  without putting them in by hand.
  \item Whether its persistent patterns can become robust particles, interact,
  and recover known symmetries and conservation laws.
  \item Whether mass, energy, geometry, and clocks can emerge from the proposed
  substrate.
  \item Whether the principle will produce a prediction unavailable to a
  complete ordinary physical account.
\end{itemize}

\section*{How the idea earns the next rung}

The next work should be less about collecting suggestive examples and more
about forcing the proposal through hard gates.

First, repeat the carrier-replacement result with different local rules,
update schedules, network shapes, and descriptions. If the effect vanishes
whenever the bookkeeping changes, it is probably an artifact.

Second, identify the smallest persistent pattern that survives replacement.
Prepare the same visible present by different histories. Then preserve, erase,
swap, scramble, or rewrite the suspected carrier while matching ordinary
energy, material, timing, and control differences. Apply the same later tests.
This asks whether the record itself does causal work or merely marks an
uncontrolled difference.

Third, build the quantum layer rather than assume it. The project must explain
how alternatives combine, why quantum probabilities have their familiar form,
and how separated systems compose. Failure here would stop the path to our
known physics even if carrier succession remained interesting on its own.

Fourth, create a true particle-like excitation. It should be possible to make
and remove it, locate it within a stated accuracy, measure how it travels,
combine or exchange it with another, scatter it from a controlled obstacle,
identify conserved features, and show that its role survives changes of local
carrier and reasonable changes of description.

Fifth, move from isolated excitations to interacting field behavior while
confronting known reasons some grid-based models fail. Only after that should
the project claim to derive mass-like or energy-like quantities. A supplied
mass knob is not an emergent mass.

Finally, make a risky physical prediction. The strongest future experiment
would state in advance where the principle and the best complete ordinary
model expect different results. If careful state enlargement always removes
the difference, if no stable carrier role survives representation changes, or
if the quantum bridge cannot be built without hiding quantum rules in the
starting machinery, the fundamental version of the program should be lowered
or abandoned.

That is not a weakness. A proposal with clear ways to fail is finally becoming
science.

\section*{Where this leaves us}

The original question survives the review in a sharper form. We do not need to
imagine information detached from the physical world. We need to understand
how an interaction builds a role, how that role can move through changing
physical realizations, and how its presence changes the futures available
under a stated test.

Science already shows that footprints can become records, records can guide
machines and organisms, and organized many-part systems can behave like new
objects. IOTA has demonstrated a small but genuine beginning: under frozen
rules, a relational bridge role can be lost locally and restored through a
causally necessary successor. The remaining distance is enormous. Quantum
structure, particles, fields, mass, energy, geometry, and a distinguishing
prediction are still ahead.

The right conclusion is neither ``proved'' nor ``just philosophy.'' It is more
interesting than either. The idea is coherent, testable, partly constructed,
and still exposed to decisive failure. We should continue---with the evidence
labels attached, the skeptical alternatives kept in the room, and every new
rung earned.

\bibliographystyle{unsrtnat}
\bibliography{library/references}

\end{document}
